\documentclass[12pt,a4paper]{article}

% === Packages ===
\usepackage[utf8]{inputenc}
\usepackage[T1]{fontenc}
\usepackage{lmodern}
\usepackage[ngerman]{babel}
\usepackage{amsmath,amssymb,amsthm}
\usepackage{graphicx}
\usepackage{booktabs}
\usepackage{hyperref}
\usepackage[margin=2.5cm]{geometry}
\usepackage{natbib}
\usepackage{algorithm}
\usepackage{algpseudocode}
\usepackage{listings}
\usepackage{xcolor}
\usepackage{float}
\usepackage{caption}
\usepackage{subcaption}
\usepackage{tikz}
\usetikzlibrary{arrows.meta,positioning,shapes.geometric,calc}

% === Hyperref config ===
\hypersetup{
    colorlinks=true,
    linkcolor=blue!60!black,
    citecolor=green!50!black,
    urlcolor=blue!70!black,
    pdftitle={Funktionaler pathway-zentrierter medizinischer Wissensgraph mit Ausschlusslogik},
    pdfauthor={Lukas Geiger},
}

% === Listings config ===
\lstset{
    basicstyle=\ttfamily\small,
    backgroundcolor=\color{gray!8},
    frame=single,
    framerule=0.4pt,
    breaklines=true,
    captionpos=b,
}

% === Theorem environments ===
\newtheorem{definition}{Definition}
\newtheorem{theorem}{Theorem}
\newtheorem{proposition}{Proposition}

% === Title ===
\title{%
    \textbf{Funktionaler pathway-zentrierter medizinischer Wissensgraph\\
    mit Ausschlusslogik}\\[0.5em]
    \large Eine systemmedizinische Architektur zur Unterstuetzung der Differentialdiagnose
}
\author{
    Lukas Geiger\\
    Unabhaengiger Forscher, Bernau, Deutschland\\
    ORCID: \url{https://orcid.org/0009-0005-7296-1534}
}
\date{Maerz 2026 --- Konzeptpapier v0.2}

\begin{document}

\maketitle

% ================================================================
\begin{abstract}
Wir stellen den Entwurf und die prototypische Implementierung eines funktionalen pathway-zentrierten medizinischen Wissensgraphen vor, der die Differentialdiagnose unterstuetzen soll. Herkoemmliche klinische Entscheidungsunterstuetzungssysteme organisieren Wissen um Diagnosen, Symptome oder Gene als primaere Entitaeten. Im Gegensatz dazu verwendet das vorgeschlagene System \emph{biologische funktionelle Pathways} als zentrales Ordnungsprinzip, wobei Gene, Proteine, Laborwerte, anatomische Lokalisationen, Zelltypen und klinische Diagnosen als sekundaere Annotationen verknuepft werden.

Diese Architektur ermoeglicht eine neuartige Form des \emph{Ausschlussschliessens}: Wenn ein funktioneller Pathway nachweislich intakt ist (ueber normale Labormarker), koennen alle Gene, die fuer diesen Pathway essentiell und hinreichend sind, als primaere Krankheitsursachen beim aktuellen Patienten ausgeschlossen werden. Das System ist als SQLite-gestuetzter Prototyp mit einer grafischen PySide6-Oberflaeche implementiert und integriert oeffentliche Datenquellen (Reactome, Gene Ontology, UniProt, HGNC, Uberon, Cell Ontology).

Wir beschreiben das formale Ausschlussmodell -- sowohl in binaerer als auch in probabilistischer Variante -- seine Annahmen und Einschraenkungen und schlagen ein Validierungsframework vor. Der Prototyp wird an einem Haemolyse-Immundefizienz-Szenario demonstriert, das fuenf funktionelle Pathways, elf Gene und sieben Labormarker umfasst. Das System ist als Forschungswerkzeug zur Erkundung des pathway-zentrierten Paradigmas positioniert, nicht als klinisches Entscheidungsunterstuetzungssystem.
\end{abstract}

\smallskip
\noindent\textbf{Schluesselwoerter:} Wissensgraph, Differentialdiagnose, funktionelle Pathways, Ausschlusslogik, Systemmedizin, Reactome, Gene Ontology

% ================================================================
\section{Einleitung und Motivation}
\label{sec:introduction}

Die Differentialdiagnose bei komplexen Multisystemerkrankungen scheitert haeufig nicht am Mangel an Daten, sondern an der Unfaehigkeit, heterogene Evidenz systematisch zu integrieren. Ein Patient mit abnormalen Laborwerten kann Hunderte genetischer Kandidatenursachen haben; aktuelle klinische Werkzeuge praesentieren diese groesstenteils als Ranglisten, ohne die logische Struktur \emph{ausgeschlossener} Moeglichkeiten auszunutzen.

Die zentrale Beobachtung, die diese Arbeit motiviert, lautet:

\begin{quote}
\emph{Wenn ein biologischer funktioneller Pathway nachweislich intakt ist, dann koennen alle Gene, die fuer diesen Pathway sowohl notwendig als auch hinreichend sind, als primaere Ursachen einer Funktionsstoerung ausgeschlossen werden.}
\end{quote}

Dies transformiert positive Evidenz (intakter Pathway) in negative diagnostische Einschraenkungen (ausgeschlossene Gene) und kann den Kandidatenraum erheblich reduzieren, bevor die Differentialdiagnose angewendet wird.

Bestehende Systeme organisieren Wissen unterschiedlich. PhenoTips \citep{Girdea2013} und Phenomizer \citep{Kohler2009} gleichen Patientenphaenotypen ueber die Human Phenotype Ontology \citep{Robinson2008} mit Krankheitsdatenbanken ab. Diese Systeme sind hervorragend im Symptom-Krankheits-Abgleich, nutzen aber nicht die Pathway-Integritaet als Ausschlussbeleg. Pathway-Datenbanken wie Reactome \citep{Fabregat2018} und KEGG \citep{Kanehisa2023} organisieren sich um biochemische Reaktionen, sind aber nicht fuer klinisches Ausschlussschliessen konzipiert. Gene Ontology \citep{Ashburner2000} liefert feingranulare Prozessbeschreibungen, es fehlt jedoch der klinische Entscheidungsrahmen.

Das vorgeschlagene System schliesst diese Luecke, indem es funktionelle Pathways als \emph{erstrangige Entitaeten} in einem Wissensgraphen behandelt, wobei alle anderen biomedizinischen Konzepte -- Gene, anatomische Lokalisationen, Zelltypen, Labormarker und Diagnosen -- als sekundaere Annotationen angebunden werden. Diese strukturelle Entscheidung ermoeglicht eine formale Ausschlusslogik, die sowohl rechnerisch handhabbar als auch biologisch sinnvoll ist.

\subsection{Beitraege}

\begin{enumerate}
    \item Ein \textbf{Datenmodell}, das funktionelle Pathways als zentrales Ordnungsprinzip fuer klinisches Schliessen verwendet und sechs Entitaetstypen ueber semantisch typisierte Kanten verknuepft.
    \item Ein \textbf{formales Ausschlussframework} mit sowohl binaerer als auch probabilistischer Variante, das die systematische Eliminierung von Kandidatengenen auf Basis von Pathway-Integritaetsevidenz ermoeglicht.
    \item Eine \textbf{prototypische Implementierung} mit Datenaufnahme aus sechs oeffentlichen biologischen Datenbanken, einem interaktiven Graphen-Explorer und einer Diagnoseabfrage-Engine.
    \item Ein \textbf{Validierungsframework}, das Metriken, Benchmark-Anforderungen und Vergleichskriterien zur Bewertung des Ansatzes spezifiziert.
\end{enumerate}

% ================================================================
\section{Problemstellung und formale Hypothese}
\label{sec:problem}

\subsection{Formale Problemstellung}

Sei $\mathcal{P} = \{p_1, \ldots, p_n\}$ die Menge der biologischen funktionellen Pathways fuer einen gegebenen Patientenkontext. Sei $\mathcal{G}$ die Menge der Kandidatengene. Sei $\mathcal{E}(p_i) \subseteq \mathcal{G}$ die Menge der fuer Pathway~$p_i$ \emph{essentiellen} Gene. Sei $s(p_i) \in \{0, 1\}$ der beobachtete Pathway-Status: $1 = \text{intakt}$, $0 = \text{gestoert oder unbekannt}$.

\begin{definition}[Binaere Ausschlussregel]
\label{def:binary}
Wenn $s(p_i) = 1$, dann gilt $\forall\, g \in \mathcal{E}(p_i)$: Gen~$g$ ist als primaere Krankheitsursache \emph{ausgeschlossen}, sofern $g$ nur in Pathways essentiell ist, die alle intakt sind:
\begin{equation}
    \text{excluded}(g) \iff \forall\, p_i \text{ s.t. } g \in \mathcal{E}(p_i): s(p_i) = 1
\end{equation}
\end{definition}

\begin{definition}[Probabilistischer Ausschlusswert]
\label{def:probabilistic}
Fuer ein Gen~$g$ ist die Ausschlusskonfi\-denz:
\begin{equation}
    \text{excl}(g) = \prod_{i:\, g \in \mathcal{E}(p_i)} c(p_i) \cdot b(g)
\end{equation}
wobei $c(p_i) \in [0,1]$ die Pathway-Statuskonfidenz ist (intakt: $0{,}95$, unbekannt: $0{,}50$, gestoert: $0{,}05$) und $b(g) \in [0,1]$ ein Multi-Pathway-Bonusfaktor:
\begin{equation}
    b(g) = \min\!\Big(1.0,\; 0.7 + 0.1 \cdot |\{p_i : g \in \mathcal{E}(p_i) \wedge s(p_i) = 1\}|\Big)
\end{equation}
Der \emph{Verdachtswert} beruecksichtigt die Genredundanz:
\begin{equation}
    \text{susp}(g) = \big(1 - \text{excl}(g)\big) \cdot r(g)
\end{equation}
wobei $r(g) \in \{1.0, 0.85, 0.6, 0.3\}$ den Redundanzgrad des Gens abbildet (keine, gering, mittel, hoch).
\end{definition}

\subsection{Kernhypothese}

\begin{quote}
\emph{Pathway-statusbasierte Ausschlusslogik reduziert die Kandidatengenmenge bei komplexen seltenen Erkrankungen systematisch um mindestens 30\,\% gegenueber rein phaenotypbasierter Filterung, wenn sie auf einem kuratierten Benchmark-Datensatz evaluiert wird.}
\end{quote}

Dies ist eine spezifische, falsifizierbare Hypothese. Die 30\,\%-Schwelle ist als minimale klinisch bedeutsame Reduktion gesetzt; ein Nullergebnis ($<10\,\%$ Reduktion) wuerde darauf hinweisen, dass die Ausschlusslogik ueber bestehende Werkzeuge hinaus wenig diagnostischen Mehrwert bietet.

Die Kandidatenreduktionsrate (CRR) ist definiert als:
\begin{equation}
\label{eq:crr}
    \text{CRR} = 1 - \frac{|\mathcal{G}_{\text{after exclusion}}|}{|\mathcal{G}_{\text{before exclusion}}|}
\end{equation}

% ================================================================
\section{Systemarchitektur}
\label{sec:architecture}

\subsection{Datenmodell}
\label{sec:datamodel}

Die zentrale Entitaet ist der \textbf{FunctionalPathway}-Knoten, der einen diskreten biologischen Prozess mit beobachtbaren Ein- und Ausgaben repraesentiert. Abbildung~\ref{fig:schema} zeigt das vollstaendige Datenmodell.

\begin{figure}[H]
\centering
\begin{tikzpicture}[
    entity/.style={draw, rounded corners, minimum width=2.8cm, minimum height=1cm, font=\small\bfseries, align=center},
    edge_label/.style={font=\scriptsize, fill=white, inner sep=1pt},
    >=Stealth,
    node distance=2.5cm and 3cm
]
    % Central node
    \node[entity, fill=blue!20] (pathway) {Funktions-\\pfad};

    % Surrounding nodes
    \node[entity, fill=green!20, above left=2cm and 3cm of pathway] (location) {K\"orper-\\ort};
    \node[entity, fill=orange!20, above right=2cm and 3cm of pathway] (cell) {Zelltyp};
    \node[entity, fill=red!20, below left=2cm and 3cm of pathway] (gene) {Gen /\\Protein};
    \node[entity, fill=purple!20, below right=2cm and 3cm of pathway] (measurement) {Messwert};
    \node[entity, fill=yellow!30, below=3cm of pathway] (diagnosis) {Diagnose};
    \node[entity, fill=purple!10, right=2cm of measurement] (signature) {Messwert-\\Signatur};
    \node[entity, fill=teal!20, above=3.5cm of pathway] (datasource) {Daten-\\quelle};

    % Edges
    \draw[->, thick] (pathway) -- node[edge_label, above, sloped] {OCCURS\_IN} (location);
    \draw[->, thick] (pathway) -- node[edge_label, above, sloped] {EXECUTED\_BY} (cell);
    \draw[->, thick] (gene) -- node[edge_label, above, sloped] {REQUIRES / MODULATES} (pathway);
    \draw[->, thick] (measurement) -- node[edge_label, above, sloped] {MONITORED\_BY} (pathway);
    \draw[->, thick] (pathway) -- node[edge_label, right] {IMPAIRED\_IN} (diagnosis);
    \draw[->, thick] (signature) -- node[edge_label, above] {COMPOSED\_OF} (measurement);
    \draw[->, dashed] (datasource) -- node[edge_label, right, pos=0.3] {PROVIDES} (pathway);

    % Cross-edges
    \draw[->, dashed] (gene) -- node[edge_label, left] {EXPRESSED\_IN} (location);
    \draw[->, dashed] (measurement) -- node[edge_label, right] {MEASURED\_AT} (location);
\end{tikzpicture}
\caption{Entity-Relationship-Diagramm des pathway-zentrierten Wissensgraphen. Durchgezogene Pfeile zeigen primaere Beziehungen ueber den Pathway-Hub; gestrichelte Pfeile zeigen sekundaere Querverbindungen.}
\label{fig:schema}
\end{figure}

Jeder Entitaetstyp traegt domaenenspezifische Attribute:

\begin{itemize}
    \item \textbf{FunctionalPathway}: Name, Beschreibung, primaere Ausgabe, Zeitklasse (akut/chronisch), Status (intakt/gestoert/unbekannt), externe~ID (Reactome oder GO).
    \item \textbf{BodyLocation} (Uberon): Organ, Subkompartiment, Zirkulationstyp, Immunrolle.
    \item \textbf{CellType} (Cell Ontology): Linie, Reifungsstadium, Lebensdauer, Mobilitaet.
    \item \textbf{Gene/Protein} (HGNC/UniProt): Symbol, Funktionstyp (strukturell, regulatorisch, signalgebend, metabolisch), Redundanzgrad, Expressionsorte.
    \item \textbf{MeasurementType}: Name, LOINC-Code, Referenzbereich, Messort, Sensitivitaet, Spezifitaet.
    \item \textbf{ClinicalDiagnosis}: Name, ICD-10-Code, Beschreibung.
\end{itemize}

Kantentypen tragen semantische Bezeichnungen (Tabelle~\ref{tab:edges}) und bei Gen-Pathway-Kanten ein \texttt{is\_essential}-Flag, das modulatorische von essentiellen Genbeitraegen unterscheidet.

\begin{table}[H]
\centering
\caption{Semantische Kantentypen im Wissensgraphen.}
\label{tab:edges}
\begin{tabular}{lll}
\toprule
\textbf{Quelle} & \textbf{Ziel} & \textbf{Relation} \\
\midrule
Pathway & BodyLocation & \texttt{OCCURS\_IN} \\
Pathway & CellType & \texttt{EXECUTED\_BY} \\
Gene & Pathway & \texttt{REQUIRES} (essentiell) / \texttt{MODULATES} \\
Measurement & Pathway & \texttt{MONITORS\_OUTPUT\_OF} \\
Measurement & BodyLocation & \texttt{REPRESENTS\_ACTIVITY\_IN} \\
Pathway & Diagnosis & \texttt{DISRUPTION\_LEADS\_TO} \\
Gene & BodyLocation & \texttt{EXPRESSED\_IN} \\
\bottomrule
\end{tabular}
\end{table}

\subsection{Implementierung}
\label{sec:implementation}

Der Prototyp verwendet \textbf{SQLite} als Backend, wobei relationale Tabellen eine Graphstruktur ueber Assoziationstabellen simulieren (Abbildung~\ref{fig:schema_sql}). Das System umfasst vier Hauptmodule:

\begin{enumerate}
    \item \textbf{Aufnahmemodul}: Laedt sechs oeffentliche Datenquellen herunter und parst sie (Abschnitt~\ref{sec:datasources}), mit einem manifestbasierten Tracking-System zur Vermeidung redundanter Downloads.
    \item \textbf{Engine-Modul}: Implementiert die Ausschlusslogik (binaer und probabilistisch), Graphtraversierungsabfragen, Mustererkennung fuer Messwert-Signaturen und die Erzeugung menschenlesbarer Erklaerungen.
    \item \textbf{GUI-Modul}: Eine PySide6-basierte Desktop-Anwendung mit vier Tabs: Graph Explorer (interaktive Netzwerkvisualisierung ueber NetworkX), Ausschluss-Panel (Pathway-Statusverwaltung und -analyse), Diagnoseabfrage (messwertgesteuerte Hypothesengenerierung) und Datenbrowser (Rohtabelleninspektion).
    \item \textbf{Datenbankmodul}: Schemadefinition, CRUD-Hilfsfunktionen und Kantenverknuepfungsfunktionen mit Fremdschluesselconstraints.
\end{enumerate}

\begin{figure}[H]
\centering
\begin{tikzpicture}[
    table/.style={draw, rounded corners=2pt, minimum width=3cm, minimum height=0.7cm, font=\ttfamily\scriptsize, fill=gray!10},
    assoc/.style={draw, rounded corners=2pt, minimum width=2.5cm, minimum height=0.6cm, font=\ttfamily\scriptsize, fill=blue!8},
    >=Stealth,
    node distance=1.2cm
]
    \node[table] (fp) {functional\_pathways};
    \node[table, left=2.5cm of fp] (bl) {body\_locations};
    \node[table, right=2.5cm of fp] (ct) {cell\_types};
    \node[table, below left=2cm and 1.5cm of fp] (gp) {genes\_proteins};
    \node[table, below right=2cm and 1.5cm of fp] (ms) {measurements};
    \node[table, below=3.5cm of fp] (dg) {diagnoses};

    \node[assoc, above=0.8cm of fp] (pl) {pathway\_locations};
    \node[assoc, above right=0.8cm and 0cm of fp] (pc) {pathway\_cells};
    \node[assoc, left=0.3cm of fp, yshift=-1cm] (gpw) {gene\_pathways};
    \node[assoc, right=0.3cm of fp, yshift=-1cm] (mp) {measurement\_pathways};
    \node[assoc, below=1.5cm of fp] (pd) {pathway\_diagnoses};

    \draw[->] (pl) -- (fp);
    \draw[->] (pl) -- (bl);
    \draw[->] (pc) -- (fp);
    \draw[->] (pc) -- (ct);
    \draw[->] (gpw) -- (gp);
    \draw[->] (gpw) -- (fp);
    \draw[->] (mp) -- (ms);
    \draw[->] (mp) -- (fp);
    \draw[->] (pd) -- (fp);
    \draw[->] (pd) -- (dg);
\end{tikzpicture}
\caption{Relationales SQLite-Schema. Graue Kaesten sind Entitaetstabellen; blaue Kaesten sind Assoziations-(Kanten-)Tabellen mit Fremdschluesseln.}
\label{fig:schema_sql}
\end{figure}

Ein dedizierter \textbf{Graph Explorer} (NetworkX-basiert) visualisiert Teilgraphen mit lazy-geladenen Nachbarschaften. Das Graph-Layout wird in einem Hintergrund-Thread berechnet, um die GUI-Responsivitaet zu erhalten. Knoten sind nach Entitaetstyp farbcodiert; Pathway-Knoten zeigen zusaetzlich ihren Status (intakt/gestoert/unbekannt) durch Randfarben an.

\paragraph{Hinweis zur Wahl der Graphdatenbank.} Die aktuelle SQLite-Implementierung ermoeglicht schnelles Prototyping, fuehrt aber zu bekannten Einschraenkungen bei tiefer Graphtraversierung (siehe Abschnitt~\ref{sec:limitations}). Ein Produktivsystem wuerde von Neo4j oder einer vergleichbaren nativen Graphdatenbank profitieren. Der Prototyp demonstriert das konzeptuelle Modell; Leistungsvergleiche mit graphnativen Backends sind geplant.

\subsection{Datenquellen}
\label{sec:datasources}

Das System integriert sechs oeffentliche biologische Datenbanken, die jeweils einen spezifischen Entitaetstyp zum Wissensgraphen beitragen:

\begin{table}[H]
\centering
\caption{Integrierte Datenquellen und ihre Rollen.}
\label{tab:datasources}
\begin{tabular}{llll}
\toprule
\textbf{Quelle} & \textbf{Entitaetstyp} & \textbf{Format} & \textbf{Status} \\
\midrule
Reactome \citep{Fabregat2018} & Pathways, Reaktionen & TSV & Integriert \\
Gene Ontology \citep{Ashburner2000} & Biologische Prozesse & OBO, GAF & Integriert \\
UniProt \citep{UniProt2023} & Proteine, Gen-Links & TSV & Integriert \\
HGNC \citep{Braschi2019} & Gennomenklatur & TSV & Integriert \\
Uberon \citep{Mungall2012} & Anatomische Lokalisationen & OBO & Integriert \\
Cell Ontology \citep{Diehl2016} & Zelltypen & OBO & Integriert \\
KEGG \citep{Kanehisa2023} & Pathway-Karten & REST API & Geplant \\
OMIM \citep{Amberger2019} / HPO \citep{Robinson2008} & Krankheit-Gen-Links & --- & Geplant \\
\bottomrule
\end{tabular}
\end{table}

Die Datenaufnahme folgt einer definierten Reihenfolge (HGNC $\to$ Uberon $\to$ Cell Ontology $\to$ UniProt $\to$ Reactome $\to$ GO), um Fremdschluesselabhaengigkeiten zu erfuellen. Jede Quelle wird ueber eine Manifestdatei verfolgt, die Download-Zeitstempel, SHA-256-Pruefsummen und Importstatus aufzeichnet.

% ================================================================
\section{Ausschlusslogik}
\label{sec:exclusion}

\subsection{Binaerer Ausschluss}

Die binaere Ausschluss-Engine (Definition~\ref{def:binary}) iteriert ueber alle Gene, die in mindestens einem Pathway als essentiell markiert sind. Fuer jedes Gen werden alle Pathways abgerufen, in denen das Gen essentiell ist, und es wird geprueft, ob \emph{alle} dieser Pathways den Status "`intakt"' haben. Ist dies der Fall, wird das Gen ausgeschlossen; andernfalls bleibt es ein Kandidat.

\begin{algorithm}[H]
\caption{Binaere Ausschlussanalyse}
\label{alg:binary}
\begin{algorithmic}[1]
\Require Menge von Pathways $\mathcal{P}$ mit Status $s(p_i)$, essentielle Gen-Zuordnung $\mathcal{E}$
\Ensure Mengen $\mathcal{G}_{\text{excluded}}$, $\mathcal{G}_{\text{candidate}}$
\State $\mathcal{I} \gets \{p_i \in \mathcal{P} : s(p_i) = \text{intakt}\}$
\ForAll{$g \in \bigcup_{p \in \mathcal{P}} \mathcal{E}(p)$} \Comment{Alle essentiellen Gene}
    \State $P_g \gets \{p \in \mathcal{P} : g \in \mathcal{E}(p)\}$ \Comment{Pathways, die $g$ benoetigen}
    \If{$P_g \subseteq \mathcal{I}$ \textbf{und} $P_g \neq \emptyset$}
        \State $\mathcal{G}_{\text{excluded}} \gets \mathcal{G}_{\text{excluded}} \cup \{g\}$
    \Else
        \State $\mathcal{G}_{\text{candidate}} \gets \mathcal{G}_{\text{candidate}} \cup \{g\}$
    \EndIf
\EndFor
\end{algorithmic}
\end{algorithm}

Die Konfidenz des Ausschlusses haengt von der Anzahl intakter Pathways ab: Ein Gen, das ueber mehrere unabhaengige intakte Pathways ausgeschlossen wird, hat staerkere Evidenz als eines, das ueber einen einzelnen Pathway ausgeschlossen wird.

\subsection{Probabilistischer Ausschluss}

Die probabilistische Engine (Definition~\ref{def:probabilistic}) erweitert das binaere Modell durch Zuweisung kontinuierlicher Konfidenzwerte. Drei Faktoren tragen bei:

\begin{enumerate}
    \item \textbf{Pathway-Evidenz}: Jeder Pathway-Status wird auf einen Konfidenzwert abgebildet (intakt: $0{,}95$, unbekannt: $0{,}50$, gestoert: $0{,}05$). Die kombinierte Pathway-Konfidenz fuer ein Gen ist das Produkt ueber alle seine essentiellen Pathways.
    \item \textbf{Multi-Pathway-Bonus}: Gene, die in mehreren intakten Pathways essentiell sind, erhalten einen Bonusfaktor, der die staerkere Evidenz aus unabhaengigen Quellen widerspiegelt.
    \item \textbf{Redundanzgewichtung}: Der Verdachtswert wird durch den Redundanzgrad des Gens moduliert. Gene ohne funktionelle Absicherung (Redundanz = "`keine"') erhalten volles Verdachtsgewicht ($r = 1{,}0$); hoch redundante Gene erhalten reduziertes Gewicht ($r = 0{,}3$).
\end{enumerate}

Gene werden nach Ausschlusskonfidenz kategorisiert: hoch ($\geq 0{,}8$), moderat ($0{,}3$--$0{,}8$) oder verdaechtig ($< 0{,}3$). Diese dreiteilige Klassifikation bietet Klinikern eine priorisierte Rangliste anstelle einer binaeren Entscheidung.

\subsection{Mustererkennung}

Zusaetzlich zum Ausschlussschliessen implementiert das System eine \textbf{Messwert-Signaturerkennung}. Eine Signatur ist ein vordefiniertes Muster abnormaler Laborwerte (z.\,B. "`Haemolyse-Signatur"': niedriges Haptoglobin + hohes LDH + hohes indirektes Bilirubin + hohe Retikulozyten + niedriges Haemoglobin). Jede Komponente traegt ein Gewicht, das ihre diagnostische Bedeutung widerspiegelt.

Bei einer Menge abnormaler Messwerte berechnet das System einen gewichteten Uebereinstimmungswert gegen alle definierten Signaturen:
\begin{equation}
    \text{match}(S) = \frac{\sum_{c \in S_{\text{matched}}} w_c}{\sum_{c \in S_{\text{all}}} w_c}
\end{equation}
Signaturen mit $\text{match} \geq 0{,}3$ werden gemeldet, was eine schnelle Identifikation bekannter klinischer Muster ermoeglicht.

% ================================================================
\section{Demonstration: Haemolyse-Immundefizienz-Szenario}
\label{sec:demo}

Der Prototyp enthaelt einen kuratierten Seed-Datensatz, der die Ueberschneidung von hereditaerer Sphaerozy\-tose, autoimmunhaemolytischer Anaemie und Agammaglobulinaemie modelliert. Dieses Szenario wurde gewaehlt, weil es mehrere interagierende Pathways mit gemeinsamen anatomischen Lokalisationen (Milz, Knochenmark) umfasst und die Ausschlusslogik ueber unabhaengige funktionelle Domaenen hinweg demonstriert.

\subsection{Funktionelle Pathways}

Fuenf Pathways werden modelliert (Tabelle~\ref{tab:pathways}), die zwei klinische Domaenen (Haematologie und Immunologie) mit unterschiedlichen zeitlichen Charakteristiken umfassen.

\begin{table}[H]
\centering
\caption{Funktionelle Pathways im Demonstrationsszenario.}
\label{tab:pathways}
\begin{tabular}{lll}
\toprule
\textbf{Pathway} & \textbf{Ausgabe} & \textbf{Zeitlich} \\
\midrule
Erythrozytenmembran-Integritaet & Stabile bikonkave Erythrozyten & Chronisch \\
Erythrozytenabbau (Haemolyse) & Freies Hb $\to$ Bilirubin & Akut \\
Milzfiltration & Entfernung defekter Zellen & Akut \\
Fruehe B-Zell-Differenzierung & Funktionelle Prae-B-Zellen & Chronisch \\
Antikoerperproduktion & IgG, IgM, IgA & Beides \\
\bottomrule
\end{tabular}
\end{table}

\subsection{Gen-Pathway-Essentialitaet}

Elf Gene sind mit diesen Pathways verknuepft, mit expliziten Essentialitaetsannotationen (Tabelle~\ref{tab:genes}). Das Essentialitaetsflag ist das entscheidende Attribut, das die Ausschlusslogik ermoeglicht: Nur essentielle Gene koennen ausgeschlossen werden, wenn ihr Pathway intakt ist.

\begin{table}[H]
\centering
\caption{Gene und ihre Pathway-Essentialitaet im Demonstrationsszenario.}
\label{tab:genes}
\begin{tabular}{llll}
\toprule
\textbf{Gen} & \textbf{Pathway} & \textbf{Essentiell} & \textbf{Redundanz} \\
\midrule
ANK1 & Membranintegritaet & Ja & Keine \\
SLC4A1 & Membranintegritaet & Ja & Keine \\
SPTA1 & Membranintegritaet & Ja & Gering \\
SPTB & Membranintegritaet & Ja & Keine \\
EPB42 & Membranintegritaet & Ja & Gering \\
HMOX1 & Haemolyse & Ja & Gering \\
HP & Haemolyse & Nein & Gering \\
BLNK & B-Zell-Differenzierung & Ja & Keine \\
CD19 & B-Zell-Differenzierung & Ja & Keine \\
PAX5 & B-Zell-Differenzierung & Ja & Keine \\
NOTCH2 & Milzfiltration, B-Zell-Diff. & Ja / Nein & Gering \\
\bottomrule
\end{tabular}
\end{table}

\subsection{Ausschlussszenario}

Betrachten wir einen Patienten mit folgender Praesentation:
\begin{itemize}
    \item Niedriges Haemoglobin (Anaemie)
    \item Niedriges Haptoglobin (Haemolyse-Marker)
    \item Hohes LDH, hohes indirektes Bilirubin (Haemolyse)
    \item Normale B-Zell-Zahl (CD19+)
    \item Normale IgG-Spiegel
\end{itemize}

Das System leitet ab:
\begin{enumerate}
    \item Normale B-Zell-Zahl $\Rightarrow$ B-Zell-Differenzierungs-Pathway \emph{intakt} $\Rightarrow$ \textbf{Ausschluss} von BLNK, CD19, PAX5.
    \item Normales IgG $\Rightarrow$ Antikoerperproduktions-Pathway \emph{intakt}.
    \item Haemolyse-Marker abnormal $\Rightarrow$ Haemolyse-Pathway \emph{gestoert}.
    \item NOTCH2 ist essentiell in der Milzfiltration (Status unbekannt) und nicht-essentiell in der B-Zell-Differenzierung $\Rightarrow$ \textbf{bleibt Kandidat}.
    \item Membrangene (ANK1, SLC4A1, SPTA1, SPTB, EPB42) sind nur im Membranintegritaets-Pathway essentiell (Status unbekannt/gestoert) $\Rightarrow$ \textbf{bleiben Kandidaten}.
\end{enumerate}

Ergebnis: 3 von 10 essentiellen Genen ausgeschlossen (CRR = 30\,\%), was der Hypothesenschwelle entspricht. Die verbleibenden Kandidaten umfassen korrekt die membranstrukturellen Gene, die mit hereditaerer Sphaero\-zytose assoziiert sind.

% ================================================================
\section{Annahmen und Einschraenkungen}
\label{sec:limitations}

\subsection{Unabhaengigkeitsannahme}

Das binaere Ausschlussmodell nimmt an, dass Pathways \emph{funktionell unabhaengig} sind: Ein intakter Pathway~$p_i$ schliesst seine essentiellen Gene unabhaengig vom Status anderer Pathways aus. Diese Annahme ist \textbf{bekanntermasssen falsch} in vielen biologischen Kontexten:

\begin{itemize}
    \item Viele Gene sind in mehreren Pathways beteiligt (\emph{Pleiotropie}). Ein ueber Pathway~$p_i$ ausgeschlossenes Gen koennte dennoch ueber einen anderen Pathway~$p_j$, der nicht erfasst wird, ursaechlich sein. Das Modell adressiert dies, indem es verlangt, dass \emph{alle} Pathways eines Gens intakt sein muessen, bevor ein Ausschluss erfolgt.
    \item Kompensationsmechanismen koennen die gemessenen Pathway-Ausgaben aufrechterhalten, selbst wenn ein essentielles Gen teilweise dysfunktional ist.
    \item Der binaere Status $s(p_i) \in \{0, 1\}$ ist eine Vereinfachung; die reale Pathway-Aktivitaet existiert auf einem Kontinuum.
\end{itemize}

Das probabilistische Modell mildert diese Einschraenkungen teilweise durch kontinuierliche Konfidenzwerte, aber die Gewichtsparameter werden derzeit heuristisch zugewiesen.

\subsection{Definition der Essentialitaet}

"`Essentielles Gen"' ist operationell definiert als ein Gen, das in Reactome oder kuratierten Quellen als \emph{erforderlich} fuer mindestens eine Reaktion im Pathway annotiert ist. Diese Definition ist konservativ und sollte unter Verwendung kuratierter Essentialitaetsdatenbanken (z.\,B. DepMap, CRISPR-Essentialitaetsscreens) verfeinert werden.

\subsection{Datenbank-Backend}

Die SQLite-Implementierung genuegt fuer den aktuellen Prototyp ($\sim$17\,MB Datenbank, $<$2000 Pathways), fuehrt aber zu Einschraenkungen bei der Vollskalierung:

\begin{itemize}
    \item Tiefe Graphtraversierungen (Tiefe $>3$) erfordern mehrfache JOIN-Operationen, die schlecht skalieren.
    \item Keine native Unterstuetzung fuer Kuerzeste-Wege- oder Musterabgleich-Abfragen, die in Graphdatenbanken verfuegbar sind (z.\,B. Neo4j Cypher).
    \item Gleichzeitiger Schreibzugriff ist durch SQLites Dateisperrung auf Dateiebene limitiert.
\end{itemize}

\subsection{Datenvollstaendigkeit}

Nicht alle Pathways haben eine ausreichende Reactome-Abdeckung. Verwaiste Pathways, gewebespezifischer Metabolismus und kuerzlich charakterisierte Mechanismen sind moeglicherweise nicht repraesentiert. Die Gene-Ontology-Integration erhoecht die Breite, jedoch auf Kosten der Granularitaet -- viele GO-Begriffe fuer biologische Prozesse sind zu generisch, um als aussagekraeftige funktionelle Pathways zu dienen.

\subsection{EHR-Interoperabilitaet}

Ein praktischer Einsatz ausschlussbasierter klinischer Entscheidungsunterstuetzung erfordert die Integration mit elektronischen Gesundheitsaktensystemen (EHR), um strukturierte Laborwerte des Patienten zu erhalten. Ohne interoperablen Zugang zu Echtzeit-Laborwerten ueber Standards wie HL7 FHIR \citep{HL7FHIR2023} bleibt das System auf manuell eingegebene Messwerte beschraenkt. EHR-Interoperabilitaet ist eine gut dokumentierte Herausforderung in der klinischen Informatik; ihre Loesung liegt ausserhalb des Rahmens dieses Konzeptpapiers, stellt aber eine kritische Anforderung fuer jede zukuenftige klinische Anwendung dar.

\subsection{Messredundanz und Kontextualitaet}

Die probabilistische Ausschlussformel (Definition~\ref{def:probabilistic}) nimmt an, dass Labormessungen unabhaengige Evidenz ueber den Pathway-Status liefern. Jedoch warnt das Contextuality-by-Default-(CbD-)Framework von Dzhafarov und Kujala \citep{Dzhafarov2016}, dass parallele Messungen informationelle Redundanz statt genuiner Unabhaengigkeit aufweisen koennen. Wenn zwei Laborwerte gleichzeitig abfallen, kann dies eine redundante Messung derselben Pathway-Ausgabe widerspiegeln, anstatt unabhaengige Evidenz fuer eine gemeinsame genetische Ursache (Pleiotropie) darzustellen. Eine rigorose Behandlung wuerde einen CbD-inspirierten Korrekturterm im Ausschlusswert erfordern, der die Ueberlappung des Messkontexts beruecksichtigt. Diese Verfeinerung ist fuer zukuenftige Versionen des probabilistischen Modells geplant.

% ================================================================
\section{Verwandte Arbeiten}
\label{sec:related}

\subsection{Phaenotypgesteuerte Diagnoseunterstuetzung}

Phenomizer \citep{Kohler2009} und PhenoTips \citep{Girdea2013} repraesentieren den Stand der Technik in der phaenotypgesteuerten Diagnose. Beide Systeme gleichen Patientenphaenotypen (HPO-Terme) mit Krankheitsdatenbanken ab und ordnen Kandidatendiagnosen nach semantischer Aehnlichkeit. Sie sind effektiv beim Symptom-Krankheits-Abgleich, modellieren aber keine biologischen Pathways und nutzen die Pathway-Integritaet nicht fuer Ausschlussschliessen.

\subsection{Pathway-Analyse in der Genomik}

Gene Set Enrichment Analysis (GSEA) \citep{Subramanian2005} und verwandte Werkzeuge identifizieren dysregulierte Pathways aus Expressionsdaten. Diese Methoden arbeiten auf Populationsebene und erfordern Omics-Daten, die in der klinischen Routineversorgung typischerweise nicht verfuegbar sind. Das vorgeschlagene System arbeitet mit individuellen Laborwerten des Patienten und ist daher in klinischen Standardarbeitsablaeufen anwendbar.

\subsection{Graphbasiertes klinisches Schliessen}

Wissensgraphen fuer die klinische Entscheidungsunterstuetzung wurden in verschiedenen Kontexten untersucht. Klinische Wissensgraphen aus elektronischen Gesundheitsakten \citep{Rotmensch2017} erfassen statistische Assoziationen zwischen Diagnosen, Laborwerten und Behandlungen. Grossskalige biomedizinische Wissensgraphen wie Hetionet \citep{Himmelstein2017} integrieren diverse Datenquellen fuer Drug-Repurposing und Hypothesengenerierung. Keines dieser Systeme nutzt jedoch den Pathway-Status als erstrangige Einschraenkung im Ausschlussschliessen.

\subsection{LLM-gestuetztes biomedizinisches Schliessen}

Neuere Arbeiten kombinieren grosse Sprachmodelle mit biomedizinischen Wissensgraphen. ESCARGOT \citep{Matsumoto2025} integriert LLMs mit einer dynamischen "`Graph of Thoughts"'-Architektur und biomedizinischen Wissensgraphen fuer Multihop-Schlussfolgerungsaufgaben. Waehrend ESCARGOT effektives graphbasiertes Schliessen in biomedizinischen Kontexten demonstriert, arbeitet es als Allzweck-Schlussfolgerungsagent ohne explizite Ausschlusslogik. Das vorgeschlagene System teilt die Betonung des wissensgraphgesteuerten Schliessens, unterscheidet sich aber grundlegend in seinem Ordnungsprinzip: funktionelle Pathways als erstrangige diagnostische Einschraenkungen statt allgemeiner Graphtraversierung.

HealthGenie \citep{Gao2025} demonstriert wissensgetriebene LLM-Integration fuer personalisierte Gesundheitsberatung. Sein Ansatz organisiert Wissensgraphen um Benutzerprofile und Ernaehrungsempfehlungen -- eine symptomzentrierte (oder bedarfszentrierte) Organisation. Im Gegensatz dazu organisiert das vorgeschlagene System um biologische funktionelle Pathways, was mechanistisches Ausschlussschliessen ermoeglicht, das in symptomzentrierten Architekturen nicht verfuegbar ist.

\subsection{Vergleichszusammenfassung}

\begin{table}[H]
\centering
\caption{Vergleich des vorgeschlagenen Systems mit verwandten Ansaetzen.}
\label{tab:comparison}
\begin{tabular}{lcccc}
\toprule
\textbf{System} & \textbf{Primaere Entitaet} & \textbf{Ausschluss} & \textbf{Laborwerte} & \textbf{Pathways} \\
\midrule
Phenomizer & HPO-Phaenotypen & Nein & Nein & Nein \\
PhenoTips & Patientenphaenotypen & Nein & Begrenzt & Nein \\
GSEA/KEGG & Pathways & Nein & Nur Omics & Ja \\
Hetionet & Heterogen & Nein & Nein & Teilweise \\
Klinische KGs & EHR-Daten & Nein & Ja & Nein \\
ESCARGOT & KG + LLM & Nein & Nein & Teilweise \\
HealthGenie & KG + LLM & Nein & Nein & Nein \\
\textbf{Vorgeschlagen} & \textbf{Pathways} & \textbf{Ja} & \textbf{Ja} & \textbf{Ja} \\
\bottomrule
\end{tabular}
\end{table}

% ================================================================
\section{Evaluationsframework}
\label{sec:evaluation}

\subsection{Benchmark-Datensatz}

Eine rigorose Validierung erfordert einen kuratierten Benchmark von 20--50 bestaetigten genetischen Diagnosen, bei denen:
\begin{enumerate}
    \item Laborprofile des Patienten zum Zeitpunkt der Erstvorstellung (vor der genetischen Diagnose) vorliegen.
    \item Die genetische Diagnose bestaetigt ist (Variante + Phaenotyp + funktionelle Studie).
    \item Der Pathway-Status des Patienten retrospektiv aus Labordaten rekonstruiert werden kann.
\end{enumerate}

Zielquellen umfassen Register fuer seltene Erkrankungen (EURORDIS), klinische Synopsen aus OMIM und publizierte Fallserien mit umfassenden Laboruntersuchungen.

\subsection{Primaere Metrik}

Die Kandidatenreduktionsrate (CRR, Gleichung~\eqref{eq:crr}) dient als primaeres Erfolgskriterium. Ein CRR $> 0{,}30$ unterstuetzt die Kernhypothese; ein CRR $< 0{,}10$ wuerde auf unzureichenden diagnostischen Mehrwert hinweisen.

\textbf{Sicherheitseinschraenkung}: Das bestaetigte ursaechliche Gen darf \emph{niemals} ausgeschlossen werden (Falschausschlussrate = 0). Jeder Falschausschluss im Benchmark invalidiert das Modell fuer die klinische Anwendung.

\subsection{Alternative Validierung: Synthetische Patientenpersonas}

Ein alternativer Validierungsweg adressiert den Mangel an annotierten realen Datensaetzen fuer seltene Erkrankungen. LLM-generierte synthetische Patientenpersonas koennen durch gezielte Stoerungen im Pathway-Graphen konstruiert werden -- beispielsweise durch Simulation eines ANK1-Defekts, indem der Erythrozytenmembranintegritaets-Pathway auf "`gestoert"' gesetzt wird und konsistente Laborwerte abgeleitet werden. Da das ursaechliche Gen \emph{a priori} definiert ist, kann die Ausschluss-Engine mit exakter Grundwahrheit evaluiert werden: Das System muss nicht-ursaechliche Gene ausschliessen (Messung der CRR) und darf das eingepflanzte ursaechliche Gen niemals ausschliessen (Messung der Falschausschlussrate). Dieser Ansatz ergaenzt den Realfall-Benchmark und ermoeglicht systematisches Testen ueber ein breiteres Spektrum von Krankheitsszenarien, als moeglicherweise aus Registern allein verfuegbar ist.

\subsection{Sekundaere Metriken}

\begin{itemize}
    \item \textbf{Rangverbesserung}: Position des korrekten Gens in der Rangkandidatenliste nach Ausschlussfilterung.
    \item \textbf{Abdeckung}: Prozentsatz der Benchmark-Faelle, bei denen mindestens ein Pathway anhand verfuegbarer Labordaten bewertet werden kann.
    \item \textbf{Vergleich}: CRR gegenueber standardmaessiger HPO-only-Filterung auf dem gleichen Fallsatz.
\end{itemize}

% ================================================================
\section{Diskussion}
\label{sec:discussion}

Der pathway-zentrierte Ansatz adressiert eine strukturelle Luecke in der aktuellen klinischen Entscheidungsunterstuetzung: die systematische Nutzung \emph{negativer Evidenz} (was nachweislich funktioniert), um den diagnostischen Raum einzuschraenken. Dies ist konzeptionell analog zur industriellen Fehlerdiagnose, bei der funktionierende Subsysteme ganze Aeste des Fehlerbaums eliminieren.

Mehrere Designentscheidungen verdienen Diskussion:

\paragraph{Warum Pathways als Zentrum, nicht Gene oder Diagnosen?} Gene sind Ursachen; Diagnosen sind Folgen. Keines von beiden bietet die mechanistische Zwischenschicht, die fuer Ausschlussschliessen benoetigt wird. Pathways repraesentieren die funktionellen Einheiten, die genetische Ursachen mit klinischen Auswirkungen verbinden, und sind damit der natuerliche Drehpunkt fuer sowohl Einschluss- als auch Ausschlusslogik.

\paragraph{Binaerer vs.\ probabilistischer Ausschluss.} Das binaere Modell ist konzeptionell einfacher und leichter zu validieren (ein Gen ist entweder ausgeschlossen oder nicht), kann aber keine Unsicherheit ausdruecken. Das probabilistische Modell erfasst abgestufte Evidenz, erfordert jedoch kalibrierte Parameter. Wir empfehlen das binaere Modell fuer die initiale Validierung und das probabilistische Modell fuer die klinische Erkundung.

\paragraph{Skalierbarkeit.} Der aktuelle Prototyp bewaeltigt das Demonstrationsszenario ($\sim$50 Knoten, $\sim$80 Kanten) ohne Leistungsprobleme. Die Skalierung auf die vollstaendige Reactome-Datenbank ($>$2400 Pathways, $>$10\,000 Reaktionen) wird Leistungstests und wahrscheinlich eine Migration zu einem Graphdatenbank-Backend erfordern.

% ================================================================
\section{Fazit und zukuenftige Arbeiten}
\label{sec:conclusion}

Der funktionale pathway-zentrierte Wissensgraph repraesentiert einen strukturell eigenstaendigen Ansatz zur Differentialdiagnose-Unterstuetzung, der Pathway-Status-Evidenz als explizite diagnostische Einschraenkungen nutzt. Die prototypische Implementierung demonstriert die Machbarkeit ueber alle Kernkomponenten hinweg: Datenaufnahme aus sechs oeffentlichen Datenbanken, Graphmodellierung mit semantischen Kantentypen, binaeres und probabilistisches Ausschlussschliessen, Mustererkennung und interaktive Visualisierung.

Die Kernhypothese -- dass Pathway-Ausschluss den Kandidatengenraum um $\geq 30\,\%$ ohne Falschausschluesse reduziert -- ist testbar und stellt das primaere empirische Ziel fuer zukuenftige Arbeiten dar.

\subsection{Zentrale offene Aufgaben}

\begin{enumerate}
    \item \textbf{Benchmark-Datensatz}: Kuratierung von 20--50 Faellen aus Registern fuer seltene Erkrankungen mit Laborprofilen und bestaetigten genetischen Diagnosen.
    \item \textbf{Gewichtskalibrierung}: Entwicklung einer datengetriebenen Methode zur Schaetzung der Pathway-Konfidenzgewichte $c(p_i)$ aus Messzuverlaessigkeit und Pathway-Redundanz.
    \item \textbf{OMIM/HPO-Integration}: Verknuepfung von Krankheit-Gen-Assoziationen aus OMIM und HPO-Annotationen, um einen direkten Vergleich mit rein phaenotypbasierter Filterung zu ermoeglichen.
    \item \textbf{Graphdatenbank-Migration}: Evaluation von Neo4j oder aehnlichen Backends fuer Skalierbarkeit ueber den SQLite-Prototyp hinaus.
    \item \textbf{Vergleich mit klinischen Werkzeugen}: Benchmark gegen Phenomizer auf demselben Fallsatz zur Quantifizierung des Mehrwerts pathway-zentrierten Ausschlusses.
    \item \textbf{Temporale Modellierung}: Erweiterung des Frameworks zur Erfassung der Krankheitsprogression und zeitabhaengiger Pathway-Statusaenderungen.
\end{enumerate}

% ================================================================
\section*{Daten- und Softwareverfuegbarkeit}

Der Quellcode des Prototyps, einschliesslich Datenbankschema, Aufnahmepipelines, Ausschluss-Engines und GUI-Anwendung, ist verfuegbar unter \url{https://github.com/lukisch/system-medicine}. Das System integriert ausschliesslich oeffentliche Datenquellen; es werden keine Patientendaten verwendet oder benoetigt.

% ================================================================
\section*{KI-Nutzungserklaerung}

Bei der Erstellung dieses Manuskripts wurden KI-gestuetzte Werkzeuge
(Claude, Anthropic) unterstuetzend eingesetzt, insbesondere fuer
Literaturrecherche, Code-Entwicklung, Textstrukturierung und Formulierungshilfe.
Die inhaltliche Konzeption, wissenschaftliche Argumentation und
Verantwortung fuer den gesamten Inhalt liegen ausschliesslich beim Autor.

% ================================================================
\bibliographystyle{apalike}
\bibliography{references}

\end{document}
